\begin{abstract}
% \bodhi{Pete's version}
% While data-driven discovery has been studied in AI for many years,
% language models (LMs) offer new opportunities to extend this paradigm,
% to not only perform automated data analysis but also, for the first time,
% handle natural language (NL) task descriptions and their
% mapping to those analyses.
% To help make progress in this broader endeavor, we present DiscoveryBench,
% the first dataset for data-driven discovery where tasks are
% specified in language that may not clearly map to data or analysis technique.
% For example, a task may mention properties not directly in, but computable
% from, the data ("total income", "BMI"), may use imprecise
% terms that need interpretation ("wealthy", "foreign"),
% and may leave the analysis approach unspecified.
% Each of the xxx tasks in DiscoveryBench are realistic and challenging,
% with data and hypotheses drawn from prior published papers so as to
% approximately reproduce the challenges those paper authors faced.
% We supplement this with
% xxx synthetic tasks, to provide additional training signal,
% and also present a new, faceted-based method for evaluating a system's
% performance on the dataset.
% Finally, we evaluate several baseline LLM-based discovery systems
% on DiscoveryBench. We find the best only scores ??\%, 
% illustrating the challenges in tackling this broader view
% of data-driven discovery, and the potential valuable role
% DiscoveryBench can play in highlighting and driving solutions to
% such challenges.


% \bodhi{Bodhi's version}
% Data-driven scientific discovery has been a paradigm of interest for many years. It encompasses the search and verification of hypotheses purely from a set of provided datasets without the need for additional data collection or physical experiments. 
% Recently, the abilities of large language models (LLMs) in code generation, function calling, and data analysis offer new opportunities to autonomously discover data-driven hypotheses.
% \ashish{shortened the opening (that would make great intro text!) and changed the pitch to start off with LLM abilities; hoping this will help highlight the relevance to LLM modeling folks}
% Can the rapidly advancing code generation, function calling, and data analysis abilities of large language models (LLMs) help automate and accelerate data-driven scientific discovery, i.e., the search and verification of hypotheses purely from a set of provided datasets?
% To help with this endeavor, we present \discoverybench{}, the first comprehensive resource formalizing the underlying multi-step hypothesis search and verification process, and designed to systematically evaluate current LLMs and build new ones for this practical task. % of data-driven discovery. 
% The 262 tasks in \discoverybench{} span 6 diverse scientific domains such as sociology, engineering, and biology. Each task is realistic and challenging and is defined by a dataset, a user's goal (in natural language), and optional relevant background knowledge about the domain.
% % , is realistic and challenging, and supports a \ashish{it seems everything would support a natural language description; is the point that we have a \emph{succinct} description? or perhaps that we have a lay-person, simple English description (as opposed to a super technical description)? at the moment, it's unclear why such a description is valuable as a selling (or distinguishing) point} natural language task description.
% Importantly, all data and hypotheses are derived from prior published papers in various domains so as to approximately reproduce the challenges those researchers faced---a key distinguishing feature from prior benchmarks on equation discovery or AutoML. We supplement this collection with 1084 synthetic tasks to provide additional training signal for models. Our proposed ontology \ashish{I'm unsure what ontology this is referring to and how that supports evaluation} \bodhi{this is our proposed ontology of context, variable, and relationship. @Dhruv may want to address this.} \ashish{oh, this doesn't sounds like an `ontology' to me; at least in CS and AI, an ontology is typically more like a structured classification of a concept into its various sub-types, e.g., WordNet, ConceptNet, Cyc, etc. How about `characterization of data-driven hypotheses'?} of data-driven hypotheses allows a systematic and multi-faceted evaluation of the discovery process. We evaluate several baseline LLM-based discovery systems on \discoverybench{}, finding that even the best scores only 16\%. This illustrates the challenges in tackling this broader view of data-driven discovery and the valuable role our benchmark can play in highlighting and driving solutions to such challenges. 
% \bhavana{should we mention diversity of scientific domains?} \bodhi{added, what do we exactly say about diversity?}
% \ashish{one point that isn't coming out well is the role of / relevance to LLMs. you mention NL task description. is that it? isn't NL description of the data (column names) also important, and so is the fact that the input is an NL query and one needs to fuzzily match that NL query to the appropriate data columns based on their NL descriptions?}

% \dhruv{Draft}
Can the rapid advances in code generation, function calling, and data analysis using large language models (LLMs) help automate the search and verification of hypotheses purely from a set of provided datasets?
To evaluate this question, we present \discoverybench{}, the first comprehensive benchmark that formalizes the multi-step process of data-driven discovery.
% \tushar{split sentence here?}
The benchmark is designed to systematically assess current model capabilities in discovery tasks and provide a useful resource for improving them.
Our benchmark contains 264 tasks collected across 6 diverse domains, such as sociology and engineering, by manually deriving discovery workflows from published papers
to approximate the real-world challenges faced by researchers, where each task is defined by a dataset, its metadata, and a discovery goal in natural language.
% and, optionally, background knowledge about the domain.
We additionally provide 903 synthetic tasks to conduct controlled evaluations across task complexity.
% and also serve as training data for model fine-tuning\tushar{If you are not fine-tuning models, it might be risky to highlight this in the abstract.}. 
Furthermore, our structured formalism of data-driven discovery enables a facet-based evaluation that provides useful insights into different failure modes. 
We evaluate several popular LLM-based reasoning frameworks using both open and closed
% \ashish{what does API-based refer to here? is it the opposite of `open' or something else? seems odd}\tushar{closed?}
LLMs as baselines on \discoverybench{} and find that even the best system scores only 25\%. 
Our benchmark, thus, illustrates the challenges in autonomous data-driven discovery and serves as a valuable resource for the community to make progress.
\end{abstract}


% Can the rapid advances in code generation, function calling, and data analysis using large language models (LLMs) help automate the search and verification of hypotheses purely from a set of provided datasets? To evaluate this question, we present DiscoveryBench, the first comprehensive benchmark that formalizes the multi-step process of data-driven discovery. The benchmark is designed to systematically assess current model capabilities in discovery tasks and provide a useful resource for improving them. Our benchmark contains 264 tasks collected across 6 diverse domains, such as sociology and engineering, by manually deriving discovery workflows from published papers to approximate the real-world challenges faced by researchers, where each task is defined by a dataset, its metadata, and a discovery goal in natural language. We additionally provide 903 synthetic tasks to conduct controlled evaluations across task complexity. Furthermore, our structured formalism of data-driven discovery enables a facet-based evaluation that provides useful insights into different failure modes. We evaluate several popular LLM-based reasoning frameworks using both open and closed LLMs as baselines on DiscoveryBench and find that even the best system scores only 25\%. Our benchmark, thus, illustrates the challenges in autonomous data-driven discovery and serves as a valuable resource for the community to make progress.